\subsection{Недообучение и переобучение, определения и способы обнаружения.}

Недообучение -- ситуация, когда модель уловила не все общие закономерности и не способна достаточно точно воспроизвести
распределение, из которого создаются объекты.
Переобучение -- ситуация, когда модель не только успешно смоделировала распределение, но и включила в него шумовые
факторы (то есть переобучилась под выбросы).
Как детектировать это все?
Вспомним про train и test, будем измерять loss на них в зависимости от сложности модели
(например, число эпох или норма вектора весов в линейной регрессии).
Если и на train, и на test loss падает, то мы уловили не все основополагающие зависимости, то есть недообучились.
Если на train loss падает, а на test loss растёт, то мы уловили шумовые зависимости, чуждые распределению, то есть
переобучились

\subsection{Способы разбиения датасета}

\subsubsection{Что такое хорошее разбиение, идеальный и практический варианты}

\textbf{Хорошее разбиение} должно удовлетворять следующим критериям:

\begin{enumerate}
    \item \textbf{Репрезентативность:} каждое подмножество (train/validation/test) должно отражать распределение
    данных генеральной совокупности
    \item \textbf{Независимость:} отсутствие пересечений между подмножествами
    \item \textbf{Адекватный размер:}
    \begin{itemize}
        \item train: достаточно для обучения модели (обычно 60-80 \%)
        \item validation: достаточно для надёжной оценки качества при настройке гиперпараметров
        \item test: достаточно для финальной оценки обобщающей способности
    \end{itemize}
    \item \textbf{Сохранение структуры:} для зависимых данных (временные ряды, панельные данные) необходимо
    сохранять временной порядок
    \item \textbf{Стратификация:} сохранение распределения целевой переменной в каждом подмножестве
\end{enumerate}

\underline{Идеальный и практический варианты}

\begin{table}[h]
    \centering
    \begin{tabular}{|p{0.45\textwidth}|p{0.45\textwidth}|}
        \hline
        \textbf{Идеальный вариант}               & \textbf{Практический вариант}          \\ \hline
        Огромный объём данных ($n > 10^6$)       & Ограниченный объем данных              \\ \hline
        Train: 80\%, Val: 10\%, Test: 10\%       & Train: 60\%, Val: 20\%, Test: 20\%     \\ \hline
        Стационарное распределение               & Возможны дрейфы распределения          \\ \hline
        i.i.d.\ наблюдения                       & Возможны зависимости                   \\ \hline
        Простое случайное разбиение              & Стратифицированное/временное разбиение \\ \hline
        Много ресурсов для перекрестной проверки & Ограниченные вычислительные ресурсы    \\ \hline
    \end{tabular}
    \caption{Сравнение идеального и практического сценариев}
    \label{tab:ideal practice}
\end{table}

\subsubsection{4 основных типа разбиений}

\underline{Простое случайное разбиение (Hold-out)}

%    \begin{algorithm}[H]
%        \caption{Простое случайное разбиение}
%        \begin{algorithmic}[1]
%            \State $n \gets \text{размер датасета}$
%            \State $\text{train\_size} \gets 0.7$, $\text{val\_size} \gets 0.15$, $\text{test\_size} \gets 0.15$
%            \State $\text{indices} \gets \text{случайная перестановка}([1, 2, \dots, n])$
%            \State $\text{train\_idx} \gets \text{indices}[1:\lfloor n \cdot \text{train\_size}\rfloor]$
%            \State $\text{val\_idx} \gets \text{indices}[\lfloor n \cdot \text{train\_size}\rfloor + 1:\lfloor n \cdot (\text{train\_size} + \text{val\_size})\rfloor]$
%            \State $\text{test\_idx} \gets \text{indices}[\lfloor n \cdot (\text{train\_size} + \text{val\_size})\rfloor + 1:]$
%        \end{algorithmic}
%    \end{algorithm}

\textbf{Преимущества}
\begin{itemize}
    \item Простота реализации
    \item Быстрота выполнения.
    \item Подходит для очень больших датасетов
\end{itemize}

\textbf{Недостатки}
\begin{itemize}
    \item Неэффективное использование данных
    \item Высокая дисперсия оценки качества
    \item Риск несбалансированного разбиения
\end{itemize}

\underline{K-Fold Cross-Validation (K-блочная перекрестная проверка)}

%    \begin{algorithm}[H]
%        \caption{K-Fold Cross-Validation}
%        \begin{algorithmic}[1]
%            \State $n \gets \text{размер датасета}$, $k \gets 5$ (обычно)
%            \State $\text{indices} \gets \text{случайная перестановка}([1, 2, \dots, n])$
%            \State $\text{folds} \gets \text{разделить indices на k равных частей}$
%            \For{$i = 1$ to $k$}
%                \State $\text{val\_idx} \gets \text{folds}[i]$
%                \State $\text{train\_idx} \gets \text{объединение всех folds, кроме i-го}$
%                \State $\text{Обучить модель на train\_idx}$
%                \State $\text{Оценить на val\_idx}$
%            \EndFor
%            \State $\text{Качество} \gets \text{среднее качество по всем folds}$
%        \end{algorithmic}
%    \end{algorithm}

\textbf{Преимущества}
\begin{itemize}
    \item Эффективное использование данных
    \item Более надёжная оценка качества
    \item Меньшая дисперсия оценки
\end{itemize}

\textbf{Недостатки}
\begin{itemize}
    \item $k$ раз больше вычислений
    \item Не подходит для зависимых данных.
    \item При больших $k$ увеличивается корреляция между folds
\end{itemize}

\underline{Стратифицированная K-Fold Cross-Validation}

%    \begin{algorithm}[H]
%        \caption{Стратифицированная K-Fold}
%        \begin{algorithmic}[1]
%            \State Сохранить распределение классов в каждом fold
%            \State Для каждого класса $c$:
%            \State \quad Разделить индексы класса на $k$ частей
%            \State Объединить соответствующие части всех классов
%        \end{algorithmic}
%    \end{algorithm}

\textbf{Преимущества}
\begin{itemize}
    \item Сохраняет распределение классов
    \item Особенно важно для несбалансированных данных
    \item Более стабильная оценка
\end{itemize}

\underline{Временное разбиение (Time-based split)}

%    \begin{algorithm}[H]
%        \caption{Временное разбиение}
%        \begin{algorithmic}[1]
%            \State Отсортировать данные по времени
%            \State $\text{cutoff\_date} \gets \text{дата разделения}$
%            \State $\text{train} \gets \text{данные до cutoff\_date}$
%            \State $\text{test} \gets \text{данные после cutoff\_date}$
%            \State (Для validation можно взять последнюю часть train)
%        \end{algorithmic}
%    \end{algorithm}

\textbf{Преимущества}
\begin{itemize}
    \item Реалистичная оценка для временных данных.
    \item Учитывает временные зависимости.
    \item Избегает утечки информации из будущего
\end{itemize}

\subsection{Этапы построения модели}

Полный конвейер построения модели состоит из 8 этапов

\underline{1. Сбор и понимание данных}
\begin{itemize}
    \item Определение бизнес-цели
    \item Сбор релевантных данных
    \item Исследовательский анализ данных (EDA)
    \item Проверка гипотез о данных
\end{itemize}

\underline{2. Предобработка и очистка данных}
\begin{itemize}
    \item Обработка пропущенных значений
    \item Обработка выбросов
    \item Кодирование категориальных признаков
    \item Масштабирование числовых признаков
\end{itemize}

\underline{3. Разбиение данных}
\begin{itemize}
    \item Выбор стратегии разбиения
    \item Создание train/validation/test наборов
    \item Проверка репрезентативности разбиения
\end{itemize}

\underline{4. Инженерия признаков}
\begin{itemize}
    \item Создание новых признаков
    \item Отбор признаков
    \item Уменьшение размерности
\end{itemize}

\underline{5. Выбор и обучение модели}
\begin{itemize}
    \item Выбор семейства моделей
    \item Обучение на train данных
    \item Валидация на validation данных
\end{itemize}

\underline{6. Настройка гиперпараметров}
\begin{itemize}
    \item Grid Search / Random Search
    \item Bayesian Optimization
    \item Выбор лучшей конфигурации
\end{itemize}

\underline{7. Оценка модели}
\begin{itemize}
    \item Финальная оценка на test данных
    \item Анализ ошибок
    \item Интерпретация результатов
\end{itemize}

\underline{8. Deploy и monitoring}
\begin{itemize}
    \item Развертывание модели
    \item Создание API
    \item Мониторинг качества
    \item Периодическое переобучение
\end{itemize}

\subsubsection{Train, validation, test, мотивация каждого}

Тренировочная выборка -- часть полной выборки, на котором мы обучаем модель.

Валидационная выборка -- часть полной выборки, происходит настройка гиперпараметров.
Это мы и назовём валидированием или валидацией модели.
Изначально человек из каких-то знаний, опыта и других соображений выбирает некоторый пул гиперпараметров для
текущей задачи.
И далее по ним с каким-то шагом запускаем модель и проверяем получившиеся предсказания с помощью метрик модели.
Выбирает и фиксирует значения гиперпараметров.

Тестовая выборка -- часть полной выборки, который мы заранее отложили и проверяем значение функции потерь, а
также метрики качества модели.
\uppercase{Не надо валидироваться на тестовой выборке}.
Считайте, что проверка на тестовой выборке -- это как финальная проверка работы.
Уже с оптимизированными параметрами и гиперапараметрами запускаем модель, чтобы проверить её качество на этой
части данных.
Иногда для конечного теста, в тестовую выборку докидывают и валидационный датасет, но такой шаг является спорным.

\subsubsection{Кросс-валидация}

Кросс-валидация -- это процедура для оценки качества работы модели, которая широко применяется в машинном обучении.
Она помогает сравнить между собой различные модели и выбрать наилучшую для конкретной задачи

Для работы модели нужно определить ее гиперпараметры.
Минус подхода с разбиениями в том, что модель будет зависеть от случайного выбора куска train (а вдруг объекты
в train упорядочены по алфавиту и не все объекты попадут в обучающую выборку при определении гиперпараметров.
Модель с этими гиперпараметрами может показать плохой результат на всей обучающей выборке).
Вторая идея (кросс-валидация): часть объектов из train идёт на валидацию, часть на train, затем данный процесс
повторяется несколько раз (часть с валидацией меняется).
Мы можем измерить дисперсию ошибки для каждого набора гиперпараметров и сам loss, чтобы выбрать оптимальный набор.
Далее речь идёт о K-fold кросс-валидации, где $K$ -- число запусков кросс-валидации и одновременно то, на сколько
примерно равных кусков мы бьём train.
Тогда в рамках каждого запуска один из $K$ блоков является валидацией, а остальные — train.
При увеличении числа блоков дисперсия ответа уменьшается (закономерно, ибо объём данных увеличивается).
Аналогичная ситуация с bias -- матожидание разности между истинным ответом и выданным алгоритмом -- он уменьшается
Чем меньше количество блоков, тем быстрее это работает, но тем меньше шанс, что модель подстроится под какой-то
кусок данных

Кросс-валидацию следует делать, когда данных мало, потому что она работает лучше.
Но, когда выборки огромные, делать кросс-валидацию очень долго.
Используем Train, validation and test stage поскольку количество данных и так позволяет хорошо обучить модель.

\subsubsection{Как выбрать организацию этапов}

\begin{table}[h]
    \centering
    \begin{tabular}{|p{0.25\textwidth}|p{0.3\textwidth}|p{0.35\textwidth}|}
        \hline
        \textbf{Тип задачи}            & \textbf{Рекомендуемая стратегия}                   & \textbf{Обоснование}                        \\\hline
        Маленький датасет ($n < 1000$) & Вложенная перекрестная проверка                    & Максимальное использование данных           \\\hline
        Несбалансированные данные      & Стратифицированная CV + oversampling/undersampling & Сохранение распределения классов            \\\hline
        Временные ряды                 & Временное разбиение + expanding window             & Сохранение временной структуры              \\\hline
        Большой датасет ($n > 10^5$)   & Hold-out + простые модели                          & Скорость вычислений важнее точности         \\\hline
        Конкурс Kaggle                 & Public/private split + ensemble                    & Избегание overfitting на public leaderboard \\\hline
    \end{tabular}
    \caption{Выбор стратегии в зависимости от задачи}
    \label{tab:stage org}
\end{table}

\underline{Принципы выбора}

\begin{enumerate}
    \item \textbf{Размер данных}
    \begin{itemize}
        \item Маленькие данные: перекрестная проверка (5-10 folds)
        \item Большие данные: hold-out (80/20 или 70/30)
    \end{itemize}
    
    \item \textbf{Распределение данных}
    \begin{itemize}
        \item Сбалансированные: простая K-fold
        \item Несбалансированные: стратифицированная K-fold
    \end{itemize}
    
    \item \textbf{Зависимости в данных}
    \begin{itemize}
        \item Независимые: случайное разбиение
        \item Временные: временное разбиение
        \item Групповые: групповые K-fold
    \end{itemize}
    
    \item \textbf{Вычислительные ресурсы}
    \begin{itemize}
        \item Ограниченные: hold-out или small K
        \item Большие: K-fold с большим K или повторная CV
    \end{itemize}
    
    \item \textbf{Цель исследования}
    \begin{itemize}
        \item Быстрый прототип: hold-out
        \item Научная статья: K-fold CV
        \item Продукт: nested CV + финальный test set
    \end{itemize}
\end{enumerate}

\underline{Вложенная перекрестная проверка (Nested Cross-Validation)}

Для надёжной оценки модели с настройкой гиперпараметров

%    \begin{algorithm}[H]
%        \caption{Nested Cross-Validation}
%        \begin{algorithmic}[1]
%            \State \textbf{Внешний цикл:} Разбиваем данные на K внешних folds
%            \For{$i = 1$ to $K$}
%                \State $\text{test\_outer} \gets \text{внешний fold } i$
%                \State $\text{train\_outer} \gets \text{остальные внешние folds}$
%
%                \State \textbf{Внутренний цикл:} Настраиваем гиперпараметры на train\_outer
%                \For{каждого набора гиперпараметров}
%                    \State Выполняем K-fold CV на train\_outer
%                    \State Выбираем лучшие гиперпараметры
%                \EndFor
%
%                \State Обучаем модель с лучшими гиперпараметрами на train\_outer
%                \State Оцениваем на test\_outer
%            \EndFor
%            \State Финальная оценка: среднее по всем внешним folds
%        \end{algorithmic}
%    \end{algorithm}

\subsection{Утечки данных}

Data leak -- утечка данных, при которой обучение происходит в том числе и на тестовых данных.
К основным типам утечек относят

%    \underline{Target Leakage (Утечка целевой переменной)}
%
%    \begin{itemize}
%        \item \textbf{Причина:} признаки, которые становятся известны только после того, как известна целевая переменная
%        \item \textbf{Пример:} прогноз оттока клиентов с признаком \("\)дата последнего платежа\("\)
%        (если клиент ушел, платежей нет)
%        \item \textbf{Решение:} тщательный анализ временных зависимостей, удаление признаков, известных только \("\)в
%        будущем \("\)
%    \end{itemize}

\underline{Train-Test Contamination (Загрязнение обучающего и тестового наборов)}

\begin{itemize}
    \item \textbf{Причина:} предобработка данных выполняется на всем датасете до разбиения
    \item \textbf{Пример:} масштабирование с вычислением mean/std по всему датасету
    \item \textbf{Решение:} все преобразования должны вычисляться только на train данных и применяться к test
\end{itemize}

\underline{3. Time Leakage (Временная утечка)}

\begin{itemize}
    \item \textbf{Причина:} использование информации из будущего для предсказания прошлого
    \item \textbf{Пример:} для предсказания курса акций используются данные, которые стали известны позже
    \item \textbf{Решение:} строгое временное разбиение, использование только исторических данных
\end{itemize}

\underline{4. Data Duplication (Дублирование данных)}

\begin{itemize}
    \item \textbf{Причина:} одинаковые или очень похожие объекты попадают и в train, и в test
    \item \textbf{Пример:} разные фотографии одного человека в разных наборах
    \item \textbf{Решение:} проверка на дубликаты перед разбиением, групповые разбиения
\end{itemize}

%    Обнаруживают и предотвращают утечки проверками
%    \underline{Процедура проверки}

%    \begin{algorithm}[H]
%        \caption{Процедура проверки на утечки}
%        \begin{algorithmic}[1]
%            \State Разделить данные на train/test ДО любой предобработки
%            \State Для всех преобразований:
%            \State \quad Вычислить параметры на train
%            \State \quad Применить к train и test отдельно
%            \State Использовать Pipeline в sklearn
%            \State Проверить распределения признаков в train и test
%            \State Если качество на тесте значительно лучше валидации --- подозрение на утечку
%        \end{algorithmic}
%    \end{algorithm}

\underline{Практические рекомендации}

\begin{enumerate}
    \item \textbf{Всегда разделяйте данные первым шагом.}
    \item \textbf{Используйте Pipeline.}
    \item \textbf{Проверяйте временные метки:} убедитесь, что в train нет данных с более поздними временными
    метками, чем в test.
    \item \textbf{Анализируйте признаки:} признаки с нереалистично высокой важностью могут указывать
    на утечку.
    \item \textbf{Сравнивайте распределения:} признаки в train и test должны иметь схожие распределения.
    \item \textbf{Используйте здравый смысл:} если модель предсказывает лучше, чем кажется возможным, проверьте на
    утечку
\end{enumerate}

\subsection*{Примеры утечек в реальных задачах}

\begin{table}[h]
    \centering
    \begin{tabular}{|p{0.3\textwidth}|p{0.3\textwidth}|p{0.3\textwidth}|}
        \hline
        \textbf{Задача}          & \textbf{Типичная утечка}                                       & \textbf{Решение}                              \\\hline
        Медицинская диагностика  & Пациент присутствует в train и test                            & Группировка по пациенту                       \\\hline
        Прогнозирование продаж   & Использование итоговых сумм месяца                             & Использовать только данные на момент прогноза \\\hline
        Рекомендательные системы & Использование будущих взаимодействий                           & Строгое временное разделение                  \\\hline
        Обработка текстов        & Обучение TF-IDF на всем корпусе                                & Обучение векторизатора только на train        \\\hline
        Анализ временных рядов   & Использование скользящего среднего с окном, включающим будущее & Только исторические данные \\\hline
    \end{tabular}
    \caption{Примеры утечек в различных задачах}
    \label{tab:leaks}
\end{table}

\subsection*{Метрики для обнаружения утечек}

\begin{itemize}
    \item \textbf{Разница в качестве:} если test accuracy $\gg$ validation accuracy.
    \item \textbf{Анализ признаков:} признаки, которые не должны влиять, имеют высокую важность.
    \item \textbf{Аномальная корреляция:} слишком высокая корреляция признаков с целевой переменной
    \item \textbf{Проверка на шум:} если модель хорошо предсказывает на полностью случайных данных
\end{itemize}

\subsection*{Заключение}

Правильное разбиение данных и избегание утечек --- фундаментальные аспекты построения надёжных ML-моделей.
Ключевые принципы:

\begin{enumerate}
    \item Разделяй данные как можно раньше.
    \item Используй Pipeline для инкапсуляции преобразований.
    \item Выбирай стратегию разбиения в зависимости от природы данных
    \item Всегда сохраняй отдельный тестовый набор для финальной оценки.
    \item Будь подозрительным к слишком хорошим результатам
\end{enumerate}